\documentclass[10pt]{extarticle}
\usepackage[utf8]{inputenc}
\usepackage[T1]{fontenc}
\usepackage[a4paper,margin=1.4cm]{geometry}
\usepackage{xcolor}
\usepackage{enumitem}
\usepackage{titlesec}
\usepackage{fancyhdr}
\usepackage{hyperref}

\definecolor{darkblue}{RGB}{15,45,90}
\definecolor{accentorange}{RGB}{220,120,30}

\newcommand{\impt}[1]{\textcolor{darkblue}{\textbf{#1}}}
\newcommand{\ntc}[1]{\textcolor{accentorange}{\textit{#1}}}

\titleformat{\section}{\large\bfseries\color{darkblue}}{\thesection.}{0.5em}{}
\titleformat{\subsection}{\normalsize\bfseries\color{accentorange}}{\thesubsection}{0.5em}{}

\pagestyle{fancy}
\fancyhf{}
\fancyhead[L]{BCC502 -- Metodologia Científica para Ciência da Computação}
\fancyhead[R]{UFOP / DECOM}
\fancyfoot[C]{\thepage}
\renewcommand{\headrulewidth}{0.4pt}

\title{\textcolor{darkblue}{\textbf{Prova -- Filosofia da Ciência}}\\[2mm]
\large Baseada em A. Chalmers, \textit{O Que É Ciência Afinal?}}
\author{BCC502 -- Metodologia Científica para Ciência da Computação}
\date{}

\begin{document}

\thispagestyle{fancy}

\begin{center}
   \large Exame Especial - BCC502 - Metodologia Científica para Ciência da Computação
\end{center}

\vspace{2mm}

\begin{enumerate}

\item Chalmers apresenta o indutivismo ingênuo e, em seguida, expõe o \impt{problema da indução} de Hume, além dos problemas relativos à base observacional (a tese de que \emph{a observação é falível e depende de teoria}).

\begin{enumerate}
\item Explique por que o problema da indução não é resolvido simplesmente por se exigir um grande número de observações feitas em condições variadas (as chamadas ``condições do indutivismo'').
\item Chalmers argumenta que os enunciados observacionais são \emph{falíveis} e \emph{dependentes de teoria}. Escolha um exemplo do seu próprio domínio de pesquisa (computação, aprendizado de máquina, dados) em que uma ``observação'' (uma métrica, um dado, um resultado experimental) só faz sentido à luz de pressupostos teóricos prévios. Explique.
\end{enumerate}



\item Distinga falseacionismo \impt{ingênuo} de falseacionismo \impt{sofisticado}, destacando o papel das hipóteses auxiliares e da tese de Duhem--Quine no segundo caso.
\item Popper defende que uma teoria científica deve ser \impt{falseável}, e que quanto mais falseável (mais proibições sobre o mundo), melhor a teoria. Um crítico poderia dizer: ``então basta formular hipóteses extremamente arriscadas e fáceis de refutar para fazer ciência de qualidade''. Mostre, usando os critérios de Popper sobre grau de falseabilidade, conteúdo empírico e verossimilhança, por que essa objeção é equivocada.
\item Segundo o falseacionismo sofisticado, uma nova teoria só deve ser aceita no lugar de outra se, além de resistir a testes, ela for \impt{independentemente testável} e tiver algum de seu excedente de conteúdo corroborado. Explique esse critério e discuta se ele resolve satisfatoriamente o problema de quando abandonar uma teoria frente a uma anomalia.





\item Explique a estrutura de um programa de pesquisa lakatosiano: \impt{núcleo firme} (hard core), \impt{heurística negativa}, \impt{heurística positiva} e \impt{cinturão protetor} de hipóteses auxiliares.
\item Lakatos propõe que um programa de pesquisa pode ser avaliado como \impt{progressivo} ou \impt{degenerativo}. Explique esse critério e argumente se ele é, de fato, uma solução melhor do que o falseacionismo popperiano para julgar o valor de uma linha de pesquisa ao longo do tempo, ou se apenas empurra o problema da decisão para outro nível (quando declarar um programa definitivamente degenerativo?).





\item Defina \impt{paradigma}, \impt{ciência normal} e \impt{revolução científica} segundo Kuhn. Explique o papel dos \impt{quebra-cabeças} (puzzles) e das \impt{anomalias} na dinâmica da ciência normal.
\item Um dos pontos mais controversos da tese de Kuhn é a noção de \impt{incomensurabilidade} entre paradigmas rivais. Explique o que essa tese afirma e por que ela levanta dificuldades para a ideia de progresso científico cumulativo. Chalmers aponta limitações dessa noção -- discuta pelo menos uma delas.





\item Explique o princípio do \impt{anarquismo epistemológico} (``tudo vale'') de Feyerabend e seu argumento de que nenhuma regra metodológica fixa resiste ao exame da história da ciência.
\item Chalmers é crítico da posição de Feyerabend. Apresente um argumento (do livro ou próprio, desde que fundamentado) contra a tese de que não existe nenhum critério racional geral para distinguir ciência de não-ciência.


\item Descreva brevemente seu projeto de pesquisa (problema, hipótese ou objetivo central, método). Em seguida, argumente qual referencial epistemológico -- falseacionismo popperiano, programa de pesquisa lakatosiano, ciência normal/paradigma kuhniano, ou mesmo uma leitura feyerabendiana -- melhor descreve a forma como seu projeto vem sendo conduzido. Não escolha o referencial que ``sonha em seguir'', mas o que de fato descreve sua prática: como você trata anomalias, como reage a resultados negativos, o que está protegido de revisão (núcleo firme) e o que está aberto a ajuste.


\item Considere a seguinte objeção, formulada por um colega:

\begin{quote}
``Seu projeto não é ciência falseável no sentido popperiano -- é apenas um cinturão protetor de hipóteses auxiliares sendo ajustado \emph{ad hoc} cada vez que o resultado não confirma o que você esperava. Isso é imunização da teoria contra a refutação, exatamente o que Popper condena.''
\end{quote}

Responda a essa objeção especificamente em relação ao referencial que você escolheu no item anterior. Sua defesa deve necessariamente: (i) reconhecer, com honestidade, algum ponto em que a crítica de fato se aplica ao seu projeto; (ii) argumentar, com base em Lakatos e/ou Kuhn, por que ajustes no cinturão protetor ou hipóteses auxiliares não são automaticamente um vício metodológico, desde que satisfaçam algum critério (ex.: testabilidade independente, geração de previsões novas, resolução de quebra-cabeças); e (iii) indicar um experimento, teste ou critério concreto que, se falhasse, faria você abandonar ou revisar substancialmente sua hipótese central -- mostrando que seu projeto não é irrefutável por princípio.


\end{enumerate}

\vspace{8mm}
\hrule
\vspace{2mm}

\end{document}
